\documentclass[11pt,a4paper]{article}

\usepackage[margin=2.4cm]{geometry}
\usepackage{amsmath,amssymb}
\usepackage{booktabs}
\usepackage{longtable}
\usepackage{microtype}
\usepackage[colorlinks=true,linkcolor=blue,urlcolor=blue]{hyperref}

\newcommand{\R}{\mathbb{R}}
\newcommand{\SEthree}{\mathrm{SE}(3)}
\newcommand{\rk}{\operatorname{rank}}
\newcommand{\B}{\mathbf{B}}
\newcommand{\Bs}{\widetilde{\mathbf{B}}}

\title{\bf Reward and observation: what each component is,\\ and how it is computed}
\author{Noyan Erdin Kilic}
\date{}

\begin{document}
\maketitle

\noindent
Reference for the environment's state score $\varphi$ and its \texttt{Dict} observation.
Source of truth: \texttt{Environment.compute\_state\_score},
\texttt{Environment.compute\_rigidity\_features}, \texttt{build\_dict\_obs} and the
primitives in \texttt{rigidity.py} / \texttt{network.py}.

\section{Notation and episode constants}

Agent $i$ has position $p_i\in\R^3$ and attitude $R_i\in SO(3)$; a \emph{domain}
$\mathcal D_i\in\{\R^2,\R^3,\R^2\!\times\!S^1,\R^3\!\times\!S^1,\SEthree\}$ fixes which
coordinates it can vary. A directed edge $i\to j$ means $i$ measures the bearing to $j$:
\[
p_{ij}=p_j-p_i,\qquad \hat p_{ij}=p_{ij}/\|p_{ij}\|,\qquad b_{ij}=R_i^{\!\top}\hat p_{ij},
\qquad P(x)=I-xx^{\!\top}.
\]
The extended bearing rigidity matrix is $\B=[\,D_p\bar E^{\!\top}\bar S \mid D_a\bar E_o^{\!\top}\bar P\,]\in\R^{3m\times 6n}$,
with one $3$-row block per directed edge, $S_i$ and $P_i$ the per-node translational and
rotational DOF projectors, and per edge $k=(i,j)$
\[
D_{p,k}=\tfrac{1}{\|p_{ij}\|}R_i^{\!\top}P(\hat p_{ij}),
\qquad
D_{a,k}=-R_i^{\!\top}[\hat p_{ij}]_\times .
\]

\paragraph{Length normalisation.} $\B$'s position columns carry $1/\text{length}$ and its
attitude columns are dimensionless, so any spectrum read off $\B$ moves when the formation
is rescaled. Everything spectral below is therefore read off
\[
\Bs=\B\cdot\operatorname{diag}(L\,I_{3n},\,I_{3n}),
\qquad
L=\Big(\tfrac1n\textstyle\sum_i\|p_i-\bar p\|^2\Big)^{1/2}
\]
(\texttt{scaled\_rigidity\_matrix}, \texttt{characteristic\_length}). Column scaling by a
positive scalar leaves the rank unchanged, so the rigidity verdict is unaffected.

\paragraph{Episode constants} (\texttt{compute\_episode\_constants}; poses only, edge
independent, computed once per \texttt{reset}).

\begin{center}
\begin{tabular}{lll}
\toprule
symbol & meaning & status \\
\midrule
$\rk_{\mathcal K}$ & $\rk \B$ of the complete graph on these poses & exact \\
$c_{\max}$ & most rank one edge block can carry, $\max_k \rk \B_{\mathcal K}[3k{:}3k{+}3,:]$ & exact \\
$m_{\text{req}}$ & fewest edges that could make these poses rigid & \emph{lower bound} \\
$L$ & RMS radius about the centroid & exact \\
$Z_{\mathcal K}$ & orthonormal basis of $\ker\Bs_{\mathcal K}$, the trivial variations & exact \\
\bottomrule
\end{tabular}
\end{center}

\noindent
$m_{\text{req}}$ is a subadditivity bound, not a ground truth, and is deliberately kept out
of the reward. Closed forms, validated exactly on 216 configurations for $n=4\ldots24$:
\[
\rk_{\mathcal K}=\sum_i \dim\mathcal D_i-q_t,
\qquad
q_t=\underbrace{(2\ \text{if any planar else}\ 3)}_{\text{translations}}+\underbrace{1}_{\text{scale}}+q_r,
\]
\[
q_r=\begin{cases}
0 & \text{some agent is } \R^2 \text{ or } \R^3 \text{ (no frame)}\\
3 & \text{every agent is } \SEthree\\
1 & \text{otherwise (all framed, shared } S^1 \text{ axis)}
\end{cases}
\qquad
c_k=\begin{cases}1 & \text{both ends planar}\\ 2 & \text{otherwise,}\end{cases}
\]
with $\dim\mathcal D_i \in \{2,3,3,4,6\}$ for the five domains in the order above, planar
meaning $\R^2$ or $\R^2\!\times\!S^1$, and $m_{\text{req}}$ the greedy accumulation of the
$c_k$ sorted descending until $\rk_{\mathcal K}$ is reached.

\section{Reward}

The step reward is potential based,
\[
r_t=\varphi(s_{t+1})-\varphi(s_t)\;-\;\text{\texttt{time\_penalty\_value}}\;+\;[\text{terminal bonus}],
\]
so it is the \emph{improvement} in $\varphi$, not its level. The discount must stay
$\gamma<1$: at $\gamma=1$ the return telescopes to $\varphi(s_T)-\varphi(s_0)$ and the
advantage collapses to $\approx 0$ under a mixing policy.

\subsection*{\texttt{WeightedNormalizedShapeError} (current)}
\[
\boxed{\;
\varphi=\frac{w_r\,\rk\B-w_e\,m\,c_{\max}}{\rk_{\mathcal K}}
\;+\;\kappa\cdot\underbrace{\frac{w_e\,c_{\max}}{\rk_{\mathcal K}}}_{\text{one edge}}\cdot\;
\mathbb 1[\text{IBR}]\cdot q\;}
\qquad w_r=100,\ w_e=25
\]
\[
q=\sigma\!\left(-\frac{\log_{10}\varepsilon-c_0}{s}\right),
\qquad
\varepsilon=\sqrt{\frac{\operatorname{tr}\big((\Bs^{\!\top}\Bs)^{+}\big)}{n}},
\qquad c_0=1.35,\ s=0.70 .
\]

\begin{itemize}\itemsep2pt
\item $m=\sum_{ij}A_{ij}$ is the edge count; $\rk\B$ comes from one thin SVD per step.
\item $m\,c_{\max}$ is ``the rank this many edges could at best have carried'', so both
      numerators are ranks and the score is dimensionless.
\item A maximally informative edge gains $(w_r-w_e)c_{\max}/\rk_{\mathcal K}$ and pruning a
      redundant one gains $w_e c_{\max}/\rk_{\mathcal K}$: an exchange rate of exactly $3{:}1$
      in every domain, at every $n$. The optimum is $w_r-w_e=75$.
\item $\varepsilon$ is \texttt{shape\_err}, read off the $\rk_{\mathcal K}$ nonzero
      eigenvalues of $\Bs^{\!\top}\Bs$; lower is better, so $q$ rises as $\varepsilon$ falls.
      It is a \emph{slope}, $\mathrm d(\text{error})/\mathrm d\sigma$, so its unit is
      $\mathrm{rad}^{-1}$ and it is not evaluated at $\sigma=1$: see \S\ref{sec:shape}
      for what it means and where it is valid.
\item $c_0,s$ are fixed constants, not per-episode quantities: no reference construction is
      built. \texttt{rigidity.SHAPE\_ERR\_EXPONENT} (default $0$) optionally subtracts
      $\alpha\log_{10}n$ from the sigmoid argument, which makes $\varphi$'s ceiling
      $n$-invariant. Re-derive all three with \texttt{tools/shape\_error\_scale.py}.
\item $\kappa=$ \texttt{stiffness\_kappa} scales the conditioning term. The term is
      bounded in $[0,\kappa\cdot\text{one edge})$, so $\kappa$ is an \emph{upper bound} on
      how many edges could ever be traded for conditioning, and it is a very loose one:
      see \S\ref{sec:kappa}. $\kappa=0$ reproduces the rank-only score exactly.
\item The term is gated on rigidity, so becoming rigid is never punished.
\end{itemize}

\subsection*{Earlier score types, kept so archived runs replay}
\begin{center}
\begin{tabular}{lll}
\toprule
type & conditioning term & per-episode cost \\
\midrule
\texttt{WeightedNormalized} & $\kappa\,(\text{one edge})\,\sigma\!\big(\log_{10}(\lambda/\lambda_{\text{ref}})/0.75\big)$ & 3 greedy constructions \\
\texttt{WeightedNormalizedSpectral} & same, functional $\in\{\lambda,\ \operatorname{tr},\ \log\det\}$ & 3 greedy constructions \\
\texttt{Weighted} & $20\rk\B-10m$, none & none \\
\bottomrule
\end{tabular}
\end{center}
$\lambda$ is the rigidity eigenvalue, the first nonzero eigenvalue of $\B^{\!\top}\B$. It has
no fixed scale, which is why those two need $\lambda_{\text{ref}}$, the median $\lambda$ of
\texttt{stiffness\_ref\_samples} greedy constructions on the episode's own poses. One
construction rebuilds $\B$ once per candidate edge. \texttt{Weighted} does not transfer
across domains: a rank-adding edge is worth $+30$ in $\R^3$ and $+10$ in $\R^2$.

\subsection{What \texttt{shape\_err} is}
\label{sec:shape}

Each bearing is a unit vector, so its noise lives in the tangent plane. With
$z_k=\text{normalize}(b_k+\sigma P(b_k)\varepsilon_k)$ and $\B$ the Jacobian of the
bearing map (checked against central differences in all five domains), the chain is
\[
\underbrace{\B^{\!\top}\B/\sigma^2}_{\text{Fisher}}
\;\longrightarrow\;
\underbrace{\sigma^2(\Bs^{\!\top}\Bs)^{+}}_{\text{Cram\'er--Rao covariance}}
\;\longrightarrow\;
\mathbb E\big[\|\delta\hat\chi-\delta\chi\|^2\big]=\sigma^2\,a_{\text{opt}},
\qquad a_{\text{opt}}=\operatorname{tr}\big((\Bs^{\!\top}\Bs)^{+}\big).
\]
So $a_{\text{opt}}$ is not a proxy for the error, it \emph{is} the expected squared error
per unit $\sigma^2$; ``A-optimality'' is only what the experiment-design literature calls
$\operatorname{tr}\Sigma$. Dividing by $n$ and taking the root gives the RMS per-agent
state error, position in formation radii and attitude in radians:
$\sigma\,\varepsilon$ with $\varepsilon=\sqrt{a_{\text{opt}}/n}$.

\paragraph{The pseudo-inverse is what makes it a \emph{shape} error.}
$\Bs^{\!\top}\Bs$ is singular and its kernel is exactly the trivial variation set
(Michieletto Thm.\ 1): translation, uniform scale, and the rotations the frames absorb,
of dimension $6n-\rk_{\mathcal K}$. The ${}^{+}$ inverts on the row space and returns zero
there, so the sum runs over the identifiable directions only and the unobservable ones are
never charged for. There is no separate gauge-quotient step. Measured at $n{=}6$
$\SEthree$: $6n=36$, $\rk_{\mathcal K}=29$, a $7$-column gauge basis with
$\|\Bs Z_{\mathcal K}\|=7.0\times10^{-15}$ and the seven smallest eigenvalues at $10^{-15}$.

\paragraph{$\sigma$ is the per-axis standard deviation, not the angle.}
The tangent plane has two independent components each of standard deviation $\sigma$, so
the RMS angle between the true and the noisy bearing is $\sigma\sqrt2$. Measured
$1.40$--$1.42$ over four decades. The Fisher information above is written in the same
per-axis convention, so every validation is self-consistent; what needs the $\sqrt2$ is
relating a noise level to a sensor specification. \texttt{--noise-sweep 1} applies a
per-axis $\sigma$ of one degree and an RMS angle of $1.41$ degrees.

\paragraph{Where it holds.} The whole construction is a linearisation. Measured over the
ratio of the recovered error to the predicted one, it is exact while the predicted
\emph{relative} error $\sigma\varepsilon$ stays below roughly $0.2$, and degrades above:

\begin{center}
\begin{tabular}{lrrrrr}
\toprule
network & $\varepsilon$ & $0.06^\circ$ & $0.57^\circ$ & $1.72^\circ$ & $5.73^\circ$ \\
\midrule
$n{=}6$ $\R^3$   &  3.6 & 1.02 & 1.02 & 0.99 & 1.01 \\
$n{=}6$ $\SEthree$ & 6.2 & 1.00 & 1.01 & 1.07 & 0.96 \\
$n{=}6$ $\R^2$   &  5.6 & 0.96 & 0.96 & 0.96 & 0.95 \\
$n{=}8$ $\R^3$   & 29.4 & 1.06 & 0.80 & 0.61 & 0.18 \\
$n{=}10$ mixed   & 22.5 & 0.95 & 1.04 & 0.87 & 0.41 \\
$n{=}10$ mixed, ill conditioned & 693.7 & 0.12 & 0.03 & 0.02 & 0.01 \\
\bottomrule
\end{tabular}
\end{center}

\noindent
The rows break at a common $\sigma\varepsilon$, not at a common $\sigma$. For a typical
$\varepsilon$ of $5$ to $30$ that is $\sigma$ below about $0.4^\circ$ to $2^\circ$. Quote
it at an operating point rather than at one radian: $\varepsilon=8$ means one degree of
per-axis bearing noise displaces the shape by about $14\%$ of its own size.

\subsection{What $\kappa$ actually trades}
\label{sec:kappa}

Over rigid graphs the score collapses to
\[
\varphi=100-\text{one edge}\cdot\big(m-\kappa q\big),
\]
so the $\varphi$-optimal graph at a given $\kappa$ is the one minimising $m-\kappa q$, and
the optimum moves from $m$ to $m+1$ edges exactly when
\[
\kappa\,\big(q^\ast(m+1)-q^\ast(m)\big)>1,
\qquad\text{i.e.}\qquad \kappa>1/\Delta q .
\]
Because $q\in(0,1)$, $\kappa$ bounds the trade at $\kappa$ edges. The realised trade is
$\kappa\,\Delta q$, and $\Delta q$ is small.

Measured by enumerating $(m,q)$ for every graph and reading the $\kappa$ curve off it
exactly: exhaustive at $n=4,5$; at $n\ge6$ the best graph per edge count is approximated
by random restarts plus single-edge-swap hill climbing, which lower bounds $q^\ast$ and so
upper bounds the $\kappa$ threshold. 46 networks over $n\in\{4,5,6,8,10\}$ and six
compositions:

\begin{center}
\begin{tabular}{lccc}
\toprule
 & range over 46 networks & typical \\
\midrule
$\Delta q$ per extra edge & $0.005$ -- $0.041$ & $0.015$ \\
$\Delta\log_{10}\varepsilon$ per extra edge & $0.027$ -- $0.187$ decades & $0.07$ \\
$\kappa$ needed to buy \emph{one} edge & $24$ -- $208$ & $\approx 65$ \\
extra edges bought at $\kappa\le20$ & \multicolumn{2}{c}{$0$, in every network} \\
extra edges bought at $\kappa=100$ & \multicolumn{2}{c}{$0$ -- $3$} \\
\bottomrule
\end{tabular}
\end{center}

\noindent
So at the $\kappa=2$ the configs use, the optimum spends \textbf{no} extra edges at all.
Part of the gap is that $q(m_{\text{req}})$ sits at $0.76$--$0.85$ rather than at the
sigmoid's centre, since $c_0$ was calibrated on greedy constructions and an exhaustive
optimum is better conditioned than one; a perfectly centred sigmoid would raise
$\mathrm dq/\mathrm dg$ from $0.23$ to at most $0.357$ and so cut the thresholds by about
$1.6\times$, leaving them an order of magnitude above $2$.

What $\kappa$ does instead is break the tie that $\kappa=0$ leaves. At $\kappa=0$ every
rigid graph on $m_{\text{req}}$ edges scores identically, and those graphs are not
equivalent at all. Exhaustively, over all rigid graphs at exactly $m_{\text{req}}$ edges:

\begin{center}
\begin{tabular}{lrrr}
\toprule
configuration & rigid graphs & worst/best $\varepsilon$ & $q$ range \\
\midrule
$n{=}4$ $\R^3$        &    48 &     $3.1\times$ & $0.686$--$0.814$ \\
$n{=}4$ $\R^2$        &   192 &     $8.4\times$ & $0.596$--$0.847$ \\
$n{=}4$ $\SEthree$    &   108 &    $77.6\times$ & $0.201$--$0.790$ \\
$n{=}4$ $\R^3\!\times\!S^1$ & 302 & $610\times$ & $0.071$--$0.803$ \\
$n{=}5$ $\R^3$        &  4480 &     $9.3\times$ & $0.543$--$0.825$ \\
$n{=}5$ mixed         & 18648 &  $7079\times$   & $0.014$--$0.782$ \\
\bottomrule
\end{tabular}
\end{center}

\noindent
$\kappa$ is therefore a \emph{selector among equally sparse graphs} over the whole usable
range, not a sparsity-versus-conditioning exchange rate. That is consistent with the greedy
sweep, where edge count is flat to within $0.08$ of an edge from $\kappa=0$ to $4$ while
stiffness moves $20\times$.

\section{Observation (\texttt{obs\_type = "Dict"})}

Five tensors. Widths below are per node or per ordered pair; bracketed flags are the env
config keys that switch a block on. Order matches \texttt{build\_dict\_obs}'s concatenation.

\subsection*{\texttt{node\_features}, $(n,\cdot)$}
\begin{longtable}{p{2.9cm}cp{9.4cm}}
\toprule
channel & w & definition \\
\midrule
\endhead
\texttt{domain} & 5 & one-hot over $(\R^2,\R^2\!\times\!S^1,\R^3,\R^3\!\times\!S^1,\SEthree)$ \\
\texttt{degree} & 2 & $(\deg^{\text{in}}_i,\deg^{\text{out}}_i)\big/(m_{\text{req}}/n)$, i.e. relative to the mean degree of a minimal graph. Dividing by $n-1$ instead introduces a $1/n$ trend \\
\texttt{closeness} & 1 & $\frac{r_i}{\sum_j d_{ij}}\cdot\frac{r_i}{n-1}$, $r_i$ reachable nodes, $d$ from Floyd--Warshall \ [\texttt{graph\_features}] \\
\texttt{eigenvector} & 1 & power iteration on $A^{\!\top}$, unit norm \ [\texttt{graph\_features}] \\
\texttt{node\_between} & 1 & Brandes betweenness \ [\texttt{graph\_features}] \\
\texttt{rigidity\_glob} & 3 & $\big(\tfrac{\rk_{\mathcal K}-\rk\B}{\rk_{\mathcal K}},\ \tfrac{m}{m_{\text{req}}},\ \mathbb 1[\text{IBR}]\big)$, tiled identically on every node \ [\texttt{rigidity\_global}] \\
\texttt{quality} & 1 & the same $q$ the reward uses, tiled; $0$ while flexible \ [\texttt{rigidity\_quality}] \\
\texttt{node\_freedom} & 1 & $\|F_i\|$ divided by its RMS over nodes, where $F$ is an orthonormal basis of $\ker\Bs\ominus\ker\Bs_{\mathcal K}$ (the non-trivial flex) and $F_i$ is node $i$'s $6$ rows. Zero once rigid \ [\texttt{rigidity\_flex}] \\
\texttt{node\_slack} & 2 & $\big(\|v_i^{\text{pos}}\|,\|v_i^{\text{att}}\|\big)$, each divided by its own mean over nodes, with $v$ the eigenvector of $\Bs^{\!\top}\Bs$ at the rigidity eigenvalue. Zero while flexible \ [\texttt{rigidity\_stiffness}] \\
\bottomrule
\end{longtable}

\subsection*{\texttt{edge\_features}, $(n,n,\cdot)$ over \emph{all} ordered pairs}
\begin{longtable}{p{2.9cm}cp{9.4cm}}
\toprule
channel & w & definition \\
\midrule
\endhead
\texttt{bearings} & 3 & $b_{ij}=R_i^{\!\top}\hat p_{ij}$ for every ordered pair, edge or not. \texttt{include\_candidate\_bearings=False} reverts to existing edges only, same shape \\
\texttt{edge\_exists} & 1 & $A_{ij}\in\{0,1\}$ \\
\texttt{edge\_between} & 1 & Brandes edge betweenness \ [\texttt{graph\_features}] \\
\texttt{reciprocity} & 1 & $A_{ij}\wedge A_{ji}$ \\
\texttt{common\_nbrs} & 1 & $(A^2)_{ij}\big/(m_{\text{req}}/n)$ \\
\texttt{add\_independence} & 1 & $\|b_{ij}Z\|_F/\|b_{ij}\|_F\in[0,1]$, $Z$ an orthonormal basis of $\ker\Bs$. Zero exactly on the pairs that would add no rank; zero everywhere once rigid \ [\texttt{rigidity\_flex}] \\
\texttt{pair\_max\_rank} & 1 & $c_k/c_{\max}$ from the \emph{complete} graph, so a candidate reads its own value rather than $0$. Constant in any homogeneous network \ [\texttt{rigidity\_edge}] \\
\texttt{add\_rank} & 1 & $\rk(b_{ij}Z)/c_{\max}$, the exact rank adding $i\to j$ would gain. Threshold $\texttt{add\_independence}>10^{-6}$, which sits mid-way in an eight-decade gap \ [\texttt{rigidity\_edge}] \\
\texttt{add\_stiffness} & 1 & $\|b_{ij}v\|$ divided by its own mean over pairs; first-order $\Delta\lambda\approx\|b_{ij}v\|^2$. A ranking prior, not an oracle. Zero while flexible \ [\texttt{rigidity\_stiffness}] \\
\texttt{remove\_rank} & 1 & rank lost by deleting $i\to j$, over $c_{\max}$: the number of unit eigenvalues of the leverage block $H=b(\Bs^{\!\top}\Bs)^{+}b^{\!\top}$. Exact. Zero on non-edges \ [\texttt{rigidity\_removal}] \\
\texttt{remove\_stiffness} & 1 & $1-\lambda(\Bs^{\!\top}\Bs-b^{\!\top}b)/\lambda(\Bs^{\!\top}\Bs)\in[0,1]$, exact via the rank-3 downdate; $1$ when removal breaks rigidity. Zero on non-edges \ [\texttt{rigidity\_removal}] \\
\bottomrule
\end{longtable}

\noindent
$b_{ij}$ is the $3\times 6n$ block edge $i\to j$ would contribute to $\Bs$; it is never
materialised, \texttt{candidate\_gain} expands
$b_{ij}Z=D_p(S_jZ_j-S_iZ_i)-D_aP_iZ_i$ into batched products over all pairs and reads norm
and rank off the $3\times3$ Gram matrix.

\subsection*{Remaining tensors}
\begin{center}
\begin{tabular}{lll}
\toprule
key & shape & definition \\
\midrule
\texttt{coord\_features} & $(n,3)$ & positions centred on the centroid, divided by $L$ \\
\texttt{adj} & $(n,n)$ & adjacency; this is what the model's action mask reads \\
\texttt{selection} & $(n)$ & pointer state, \texttt{SelectNodesSequentially} only \\
\bottomrule
\end{tabular}
\end{center}

\section{Invariance}

The task is defined up to a similarity transform, so the observation and $\varphi$ should be
too. Measured by transforming the network and diffing each channel, both regimes
(flexible and rigid), $\R^3$ / $\SEthree$ / mixed, $n$ up to $100$.

\begin{center}
\begin{tabular}{lccc}
\toprule
 & translation & rotation & uniform scaling \\
\midrule
$\varphi$ & exact & exact & exact \\
every rigidity channel & $\le 10^{-11}$ & $\le 10^{-11}$ & $\le 10^{-11}$ \\
\texttt{degree}, \texttt{adj}, \texttt{reciprocity}, \texttt{common\_nbrs} & exact & exact & exact \\
\texttt{bearings} & exact & \textbf{moves in $\R^2,\R^3$} & exact \\
\texttt{coord\_features} & exact & \textbf{equivariant by design} & exact \\
\bottomrule
\end{tabular}
\end{center}

\begin{itemize}\itemsep2pt
\item \texttt{bearings} rotate in the frameless domains because $R_i=I$ there, so
      \texttt{get\_bearing} returns a global-frame vector. Exact to $10^{-15}$ in
      $\R^2\!\times\!S^1$, $\R^3\!\times\!S^1$, $\SEthree$. \texttt{rotation\_augmentation}
      is the mitigation.
\item \texttt{coord\_features} rotate deliberately. GINE never receives them; EGNN reads
      them only through $\|x_i-x_j\|^2$, which is invariant.
\item Scale invariance of the three stiffness channels requires that $v$, $Z$ and the
      removal downdate all come from $\Bs$ rather than from $\B$. Mapping a null
      \emph{space} into scaled units afterwards is valid; applying the same map to the
      single eigenvector $v$ is not.
\item Channels normalised by their own mean or RMS (\texttt{node\_freedom},
      \texttt{node\_slack}, \texttt{add\_stiffness}) encode \emph{which} node or pair is
      soft, never how soft in absolute terms.
\end{itemize}

\section{Scale under a per-episode domain draw}

With \texttt{random\_domains}, the composition is redrawn every \texttt{reset}, so
$\rk_{\mathcal K}$, $c_{\max}$ and $m_{\text{req}}$ move between episodes and with them the
size of one reward step. Measured over 4000 uniform draws per $n$:

\begin{center}
\begin{tabular}{rccc}
\toprule
$n$ & one edge, median [p10--p90] & p90/p10 & $\varphi$ optimum p90/p10 \\
\midrule
8   & 1.923 [1.613--2.381] & 1.48 & 1.015 \\
16  & 0.926 [0.806--1.042] & 1.29 & 1.007 \\
32  & 0.446 [0.410--0.490] & 1.20 & 1.003 \\
64  & 0.220 [0.207--0.234] & 1.13 & 1.002 \\
100 & 0.140 [0.133--0.147] & 1.11 & 1.001 \\
\bottomrule
\end{tabular}
\end{center}

\begin{itemize}\itemsep2pt
\item The spread \emph{narrows} with $n$: $\rk_{\mathcal K}$ is a sum of $n$ i.i.d. per-agent
      DOFs, so its relative spread falls as $1/\sqrt n$. The corner-to-corner range is a
      flat $2.0\times$ at every $n$; uniform draws occupy a shrinking part of it.
\item $\varphi$'s optimum is effectively constant, which is what a value-based method needs:
      with a potential-based reward the return telescopes to $\varphi(s_T)-\varphi(s_0)$.
\item $c_{\max}=1$ requires every agent planar: $0.05\%$ of draws at $n=8$, none above.
\item Within one episode the realised rewards already span $-3$ to $+3$ times one edge.
      Measured across 120 episodes, the within-episode standard deviation exceeds the
      between-episode one by $2$--$3.8\times$, and switching from a fixed mix to a random
      one raises the between-episode figure from $0.434$ to $0.507$ at $n=10$ and not at
      all at $n=20$.
\item Across \emph{sizes} the step size does move by an order of magnitude ($1.92$ at $n=8$
      to $0.109$ at $n=128$), since one edge is $w_ec_{\max}/\rk_{\mathcal K}$ and
      $\rk_{\mathcal K}$ grows linearly in $n$.
\end{itemize}

\end{document}
